\documentclass[12pt,a4paper,oneside]{scrreprt}

\usepackage[utf8]{inputenc}
\usepackage[T1]{fontenc}
\usepackage{lmodern}
\usepackage[english]{babel}
\usepackage{amsmath,amssymb,amsfonts}
\usepackage{bm}
\usepackage{booktabs}
\usepackage{csquotes}
\usepackage{graphicx}
\usepackage{microtype}
\usepackage{siunitx}
\usepackage{subcaption}
\usepackage{xcolor}
\usepackage{hyperref}

\hypersetup{
  colorlinks=true,
  linkcolor=blue!50!black,
  citecolor=blue!50!black,
  urlcolor=blue!50!black
}

\graphicspath{{figures/}{../figures/}}

\newcommand{\SU}{\mathrm{SU}}
\newcommand{\U}{\mathrm{U}}
\newcommand{\dd}{\mathrm{d}}
\newcommand{\order}{\mathcal{O}}
\newcommand{\eps}{\varepsilon}
\newcommand{\latspacing}{a}
\newcommand{\wilsonloop}{W(R,T)}
\newcommand{\creutz}{\chi(R,T)}
\newcommand{\rzero}{r_0}
\newcommand{\stringtension}{\sigma}

\title{Investigation of the Continuum Limit in the Presence of Small Breaking of Gauge Invariance in Lattice Gauge Theories}
\author{Piet Leyendecker}
\date{\today}

\begin{document}

\maketitle

\begin{abstract}
This thesis investigates whether and how the continuum limit of a lattice gauge theory is affected by small explicit violations of gauge invariance. The study [...] numerical measurements of Wilson loops, Creutz ratios,
 the static potential, the Sommer scale, and related scaling observables. The central goal is to determine
 whether gauge-breaking effects vanish, persist, or modify the extrapolation toward the continuum limit as the lattice spacing is reduced.
\end{abstract}

\tableofcontents

\chapter*{Notation and Conventions}
\addcontentsline{toc}{chapter}{Notation and Conventions}
\begin{description}
  \item[$\latspacing$] Lattice spacing.
  \item[$\beta$] Bare lattice coupling parameter.
  \item[$\eps$] Parameter controlling the explicit breaking of gauge invariance.
  \item[$\wilsonloop$] Wilson loop with spatial extent $R$ and temporal extent $T$.
  \item[$\creutz$] Creutz ratio used to estimate the force and string tension.
  \item[$\rzero$] Sommer scale defined through $r^2F(r)=1.65$.
  \item[$\stringtension$] String tension.
\end{description}

\chapter{Introduction}
\section{Motivation}
Explain why gauge invariance is normally central in lattice gauge theory and why controlled violations are relevant. Possible motivations include numerical robustness, effective descriptions, regulator artifacts, quantum simulation errors, or imperfect implementation of gauge constraints.

\section{Research Question}
State the main question precisely:
\begin{quote}
  Does a small explicit breaking of gauge invariance disappear in the continuum limit, or does it alter the continuum extrapolation of physical observables?
\end{quote}

\section{Strategy and Main Observables}
Introduce the numerical strategy and the observables used to probe the continuum limit: Wilson loops, Creutz ratios, the static potential $V(R)$, $\rzero$, and $\stringtension$.

\chapter{Theoretical Background}
\label{chap:theoretical-background}

The aim of this chapter is to collect the theoretical ingredients needed for the
analysis of the continuum limit in the presence of small explicit violations of
gauge invariance. The discussion starts from continuum Yang--Mills theory and
then introduces the lattice formulation, Wilson loops, the static potential,
Creutz ratios, and the Sommer scale. These are the central concepts used later
to define and compare physical observables at different lattice spacings.

\section{Gauge Theories in the Continuum}
\label{sec:continuum-gauge-theory}

\subsection{Gauge Fields and Local Gauge Symmetry}
\label{subsec:gauge-fields-local-symmetry}

A non-Abelian gauge theory is built from fields that transform locally under a
Lie group. In this thesis the relevant case is the group $\SU(2)$. A continuum
gauge field can be written as
\begin{equation}
  A_\mu(x) = A_\mu^a(x) T^a ,
\end{equation}
where $T^a$ are the generators of the Lie algebra of $\SU(2)$. A common choice is
$T^a = \sigma^a/2$, with $\sigma^a$ the Pauli matrices. The gauge field defines
the covariant derivative
\begin{equation}
  D_\mu = \partial_\mu + i A_\mu ,
\end{equation}
and the associated field-strength tensor
\begin{equation}
  F_{\mu\nu}
  =
  \partial_\mu A_\nu
  -
  \partial_\nu A_\mu
  +
  i [A_\mu,A_\nu] .
\end{equation}
The commutator term distinguishes a non-Abelian gauge theory from an Abelian
one. It implies that the gauge field itself carries charge under the gauge
group and therefore self-interacts.

A local gauge transformation is specified by a group element
\begin{equation}
  \Omega(x) \in \SU(2)
\end{equation}
at every point in space-time. Under such a transformation, matter fields in the
fundamental representation transform as
\begin{equation}
  \psi(x) \mapsto \Omega(x)\psi(x),
\end{equation}
while the gauge field transforms inhomogeneously,
\begin{equation}
  A_\mu(x)
  \mapsto
  \Omega(x) A_\mu(x) \Omega^\dagger(x)
  -
  i \bigl(\partial_\mu \Omega(x)\bigr)\Omega^\dagger(x).
\end{equation}
The inhomogeneous derivative term is necessary because the transformation is
local. The field strength transforms covariantly,
\begin{equation}
  F_{\mu\nu}(x)
  \mapsto
  \Omega(x) F_{\mu\nu}(x) \Omega^\dagger(x),
\end{equation}
so traces of products of field strengths at the same point are gauge invariant.

In Euclidean space-time, the pure Yang--Mills action can be written as
\begin{equation}
  S_{\mathrm{YM}}
  =
  \frac{1}{2g^2}
  \int \dd^4 x \,
  \mathrm{tr}\,
  F_{\mu\nu}(x) F_{\mu\nu}(x),
  \label{eq:continuum-yang-mills-action}
\end{equation}
with summation over repeated Lorentz indices. This action is invariant under
local gauge transformations. The coupling $g$ controls the strength of the
interaction. For an asymptotically free gauge theory, the effective coupling
becomes small at short distances, while the long-distance theory is strongly
coupled and requires non-perturbative methods.

\subsection{Gauge-Invariant Observables}
\label{subsec:gauge-invariant-observables}

Physical observables in a gauge theory must be gauge invariant. Gauge-dependent
quantities can be useful after a gauge fixing has been imposed, but they do not
directly correspond to physical observables unless the gauge dependence is
properly controlled. For pure gauge theory, important gauge-invariant
observables include traces of closed parallel transporters, local action-density
operators, correlation functions of gauge-invariant operators, and quantities
derived from Wilson loops.

A central continuum object is the path-ordered parallel transporter along a
curve $C$,
\begin{equation}
  U(C)
  =
  \mathcal{P}
  \exp\left(
    i \int_C \dd x^\mu A_\mu(x)
  \right).
\end{equation}
For a closed curve, the trace
\begin{equation}
  W(C)
  =
  \frac{1}{N_c}
  \mathrm{Re}\,\mathrm{tr}\, U(C)
\end{equation}
is gauge invariant. On the lattice this object becomes the Wilson loop and is
one of the most direct probes of confinement and of the static quark--antiquark
potential.

\subsection{Physical Consequences of Gauge Invariance}
\label{subsec:consequences-gauge-invariance}

Gauge invariance is not only a formal redundancy of the description. It
constrains the allowed terms in the action, the form of correlation functions,
and the possible continuum limits. In the lattice formulation, preserving gauge
invariance exactly at nonzero lattice spacing is usually advantageous because it
prevents gauge-noninvariant counterterms from being generated by radiative
corrections. Wilson's lattice gauge action realizes this idea by putting compact
group-valued variables on links and by constructing the action from traces of
closed loops \cite{Wilson1974,GattringerLang2010,MontvayMuenster1994}.

The purpose of this thesis is to study what happens when this exact symmetry is
slightly violated by an explicit deformation. From a renormalization-group point
of view, the key question is whether the deformation behaves like an irrelevant
lattice artifact, so that its effects vanish as $\latspacing \to 0$, or whether
it defines a relevant perturbation that changes the continuum limit. This
question is not purely formal: small violations of gauge constraints can appear
as regulator artifacts, algorithmic imperfections, or implementation errors in
effective or quantum-simulation formulations of lattice gauge theories
\cite{Foerster1980,Bonati2021Breaking}.

\section{Lattice Gauge Theory}
\label{sec:lattice-gauge-theory}

\subsection{Discretization of Space-Time}
\label{subsec:discretization}

In lattice gauge theory, Euclidean space-time is replaced by a hypercubic
lattice
\begin{equation}
  \Lambda_{\latspacing}
  =
  \left\{
    x = \latspacing n
    \,\middle|\,
    n \in \mathbb{Z}^4
  \right\},
\end{equation}
where $\latspacing$ is the lattice spacing. The lattice spacing acts as an
ultraviolet cutoff. Momenta larger than order $1/\latspacing$ are not resolved.
The continuum limit is obtained by sending
\begin{equation}
  \latspacing \to 0
\end{equation}
while keeping physical quantities fixed.

In a numerical simulation the lattice has finite extent
\begin{equation}
  L_\mu = N_\mu \latspacing ,
\end{equation}
where $N_\mu$ is the number of lattice points in direction $\mu$. Therefore two
limits are conceptually distinct: the continuum limit $\latspacing \to 0$ and
the infinite-volume limit $L_\mu \to \infty$. In practice, one tries to choose
volumes large enough that finite-volume effects are smaller than the statistical
and systematic errors relevant for the analysis. [in simulations in this thesis I use Lµ=const for the space directions and Nµ const for time direction]

\subsection{Link Variables and Plaquettes}
\label{subsec:link-variables-plaquettes}

The fundamental degrees of freedom of Wilson lattice gauge theory are link
variables
\begin{equation}
  U_\mu(x) \in \SU(2),
\end{equation}
which live on the oriented link from $x$ to $x+\latspacing\hat\mu$. They are the
lattice analogues of parallel transporters. For a smooth continuum field one may
formally write
\begin{equation}
  U_\mu(x)
  \approx
  \exp\left[
    i \latspacing A_\mu\left(x+\frac{\latspacing}{2}\hat\mu\right)
  \right].
\end{equation}
A local lattice gauge transformation assigns a group element $\Omega(x)$ to each
site and acts on the link variables as
\begin{equation}
  U_\mu(x)
  \mapsto
  \Omega(x) U_\mu(x) \Omega^\dagger(x+\latspacing\hat\mu).
  \label{eq:lattice-gauge-transformation}
\end{equation}

The elementary closed loop on the lattice is the plaquette,
\begin{equation}
  U_{\mu\nu}(x)
  =
  U_\mu(x)
  U_\nu(x+\latspacing\hat\mu)
  U_\mu^\dagger(x+\latspacing\hat\nu)
  U_\nu^\dagger(x).
  \label{eq:plaquette-definition}
\end{equation}
Under a gauge transformation, the plaquette transforms covariantly at its base
point,
\begin{equation}
  U_{\mu\nu}(x)
  \mapsto
  \Omega(x) U_{\mu\nu}(x) \Omega^\dagger(x).
\end{equation}
Therefore $\mathrm{Re}\,\mathrm{tr}\,U_{\mu\nu}(x)$ is gauge invariant.

For smooth fields, the plaquette has the continuum expansion
\begin{equation}
  U_{\mu\nu}(x)
  =
  \exp\left[
    i \latspacing^2 F_{\mu\nu}(x)
    +
    \order(\latspacing^3)
  \right].
\end{equation}
Thus the plaquette is the lattice object from which the continuum field
strength is recovered. [but not really used in this thesis, so maybe remove this point]

\subsection{Wilson Gauge Action}
\label{subsec:wilson-gauge-action}

The standard Wilson gauge action for $\SU(N_c)$ is
\begin{equation}
  S_W[U]
  =
  \beta
  \sum_x
  \sum_{\mu<\nu}
  \left(
    1
    -
    \frac{1}{N_c}
    \mathrm{Re}\,\mathrm{tr}\,
    U_{\mu\nu}(x)
  \right),
  \label{eq:standard-wilson-action}
\end{equation}
with
\begin{equation}
  \beta = \frac{2N_c}{g_0^2}.
\end{equation}
For $\SU(2)$ this becomes
\begin{equation}
  \beta = \frac{4}{g_0^2},
\end{equation}
and
\begin{equation}
  S_W[U]
  =
  \beta
  \sum_x
  \sum_{\mu<\nu}
  \left(
    1
    -
    \frac{1}{2}
    \mathrm{Re}\,\mathrm{tr}\,
    U_{\mu\nu}(x)
  \right).
  \label{eq:wilson-action-su2}
\end{equation}
The action is built only from traces of closed loops and is therefore exactly
gauge invariant at nonzero lattice spacing.

The continuum limit of the Wilson action is obtained by inserting the plaquette
expansion. Up to an irrelevant additive constant and lattice artifacts,
\begin{equation}
  S_W[U]
  \longrightarrow
  \frac{1}{2g_0^2}
  \int \dd^4x\,
  \mathrm{tr}\,
  F_{\mu\nu}(x)F_{\mu\nu}(x).
\end{equation}
This reproduces the Euclidean Yang--Mills action. The corrections at finite
lattice spacing are cutoff effects. For the pure Wilson gauge action, the
leading gauge-invariant lattice artifacts in many on-shell quantities are of
order $\latspacing^2$, although the actual scaling observed in numerical data
can be affected by statistics, finite-volume effects, operator definitions, and
the accessible range of lattice spacings.

\subsection{Path Integral and Expectation Values}
\label{subsec:lattice-path-integral}

The lattice path integral is finite-dimensional at finite lattice spacing and
finite volume. The partition function is
\begin{equation}
  Z
  =
  \int
  \prod_{x,\mu}
  \dd U_\mu(x)
  \,
  e^{-S[U]},
  \label{eq:lattice-partition-function}
\end{equation}
where $\dd U_\mu(x)$ is the Haar measure on $\SU(2)$. The expectation value of
an observable $\mathcal{O}[U]$ is
\begin{equation}
  \langle \mathcal{O} \rangle
  =
  \frac{1}{Z}
  \int
  \prod_{x,\mu}
  \dd U_\mu(x)
  \,
  \mathcal{O}[U]
  e^{-S[U]}.
  \label{eq:lattice-expectation-value}
\end{equation}
Gauge-invariant expectation values can be computed without gauge fixing. This
is one of the practical advantages of the compact Wilson formulation.

In Monte Carlo simulations the path integral is estimated by generating an
ensemble of gauge-field configurations distributed according to $e^{-S[U]}$.
For an ensemble of $N_{\mathrm{cfg}}$ configurations $U_i$, the estimator is
\begin{equation}
  \langle \mathcal{O} \rangle
  \approx
  \frac{1}{N_{\mathrm{cfg}}}
  \sum_{i=1}^{N_{\mathrm{cfg}}}
  \mathcal{O}[U_i].
\end{equation}
Because successive configurations in a Markov chain are correlated,
autocorrelations must be considered when estimating uncertainties.

\subsection{Continuum Limit and Renormalization}
\label{subsec:continuum-limit-renormalization}

The continuum limit is not obtained by simply decreasing $\latspacing$ at fixed
bare coupling. Instead, the bare parameters have to be tuned so that physical
observables remain fixed in physical units. For pure $\SU(2)$ gauge theory with
the Wilson action, decreasing the lattice spacing corresponds to increasing
$\beta$, or equivalently decreasing the bare coupling $g_0$.

A useful way to formulate the continuum limit is in terms of a physical length
scale, such as the Sommer scale $\rzero$. The dimensionless ratio
\begin{equation}
  \frac{\rzero}{\latspacing}
\end{equation}
is measured on the lattice. Increasing $\rzero/\latspacing$ means that the same
physical length is resolved by more lattice points, i.e. the lattice spacing is
smaller in physical units. Therefore continuum extrapolations can be expressed
as fits in
\begin{equation}
  \left(\frac{\latspacing}{\rzero}\right)^2
\end{equation}
or another measured dimensionless lattice-spacing proxy.

Continuum extrapolations compare observables at fixed physical
conditions, not merely at fixed lattice indices. For observables derived from
Wilson loops this means that a loop size $R$ on a coarse ensemble and the same
integer $R$ on a fine ensemble do not represent the same physical distance.
Therefor instead compare at fixed ratios such as $r/\rzero$ or fixed
dimensionless distances, a reference scale is used.
[This point is emphasized in recent]
work on controlled continuum extrapolations of Wilson-loop observables, where
Creutz ratios are used as finite continuum quantities and fitted at fixed
physical distance \cite{Risch2024SmearingContinuum}.

The renormalization-group interpretation is central for the present thesis.
Different lattice actions can lead to the same continuum theory if they differ
only by irrelevant operators. In that case they lie in the same universality
class. By contrast, a relevant deformation changes the long-distance theory
unless its coefficient is tuned to zero. [A gauge-breaking term must therefore be
classified by its effect on long-distance, gauge-invariant observables as
$\latspacing \to 0$. (not written well yet)] The numerical analysis in later chapters tests this
question directly.

\section{Static Potential and Confinement}
\label{sec:static-potential-confinement}

\subsection{Wilson Loops}
\label{subsec:wilson-loops}

A Wilson loop is the trace of the ordered product of link variables around a
closed contour. For a rectangular loop in the $\mu\nu$-plane with side lengths
$L_\mu$ and $L_\nu$ in lattice units, the code measures the normalized real
trace
\begin{equation}
  W_{\mu\nu}(L_\mu,L_\nu)
  =
  \frac{1}{N_c V}
  \sum_x
  \operatorname{Re}\operatorname{tr}
  \left[
    U_\mu(x;L_\mu)\,
    U_\nu(x+L_\mu\hat\mu;L_\nu)\,
    U_\mu^\dagger(x+L_\nu\hat\nu;L_\mu)\,
    U_\nu^\dagger(x;L_\nu)
  \right],
  \label{eq:rectangular-wilson-loop-code}
\end{equation}
where $V$ is the lattice volume and
\begin{equation}
  U_\mu(x;L_\mu)
  =
  U_\mu(x)
  U_\mu(x+\hat\mu)
  \cdots
  U_\mu(x+(L_\mu-1)\hat\mu)
\end{equation}
denotes a straight Wilson line of length $L_\mu$ in direction $\mu$. The factor
$1/N_c$ normalizes the trace, while the factor $1/V$ averages over all possible
base points of the loop.

For temporal Wilson loops, the last lattice direction is used as the Euclidean
time direction,
\begin{equation}
  t = N_d - 1.
\end{equation}
The code computes rectangular loops in each spatial-time plane and averages over
the spatial directions:
\begin{equation}
  W(R,T)
  =
  \frac{1}{N_d-1}
  \sum_{\mu=0}^{N_d-2}
  W_{\mu,t}(R,T).
  \label{eq:temporal-wilson-loop-spatial-average}
\end{equation}
For the $\SU(2)$ simulations considered in this thesis, $N_c=2$, so the stored
temporal Wilson loop is
\begin{equation}
  W(R,T)
  =
  \frac{1}{2V(N_d-1)}
  \sum_{\mu=0}^{N_d-2}
  \sum_x
  \operatorname{Re}\operatorname{tr}
  W_{\mu,t}(x;R,T).
  \label{eq:temporal-wilson-loop-su2-code}
\end{equation}
Thus the values written to the output file are already normalized by the color
factor, averaged over the full lattice volume, and averaged over spatial
directions.

In the implementation, the rectangular loop is constructed from two paths that
connect opposite corners of the rectangle. The first path is
\begin{equation}
  l_\mu
  =
  U_\mu(x;L_\mu)\,
  U_\nu(x+L_\mu\hat\mu;L_\nu),
\end{equation}
and the second path is
\begin{equation}
  l_\nu
  =
  U_\nu(x;L_\nu)\,
  U_\mu(x+L_\nu\hat\nu;L_\mu).
\end{equation}
The measured loop contribution is then
\begin{equation}
  \operatorname{Re}\operatorname{tr}
  \left(
    l_\mu l_\nu^\dagger
  \right).
\end{equation}
This is equivalent to the trace of the ordered product around the rectangular
closed contour.

The temporal Wilson-loop output file has the structure
\begin{verbatim}
# step, L, T, W_temp
\end{verbatim}
where \texttt{L} is the spatial extent $R$, \texttt{T} is the temporal extent,
and \texttt{W\_temp} is the averaged, normalized value defined in
Eq.~\eqref{eq:temporal-wilson-loop-spatial-average}.

\subsection{Effective Potential Extraction}
\label{subsec:effective-potential-extraction}

Temporal Wilson loops are used to extract the static quark--antiquark potential.
For large Euclidean time extent $T$, the Wilson loop has the spectral form
\begin{equation}
  W(R,T)
  =
  \sum_n c_n(R) e^{-a E_n(R) T},
  \label{eq:wilson-loop-spectral-decomposition}
\end{equation}
where $E_0(R)=V(R)$ is the static ground-state potential and the higher
$E_n(R)$ describe excited states with the same static sources. At sufficiently
large $T$, the ground state dominates:
\begin{equation}
  W(R,T)
  \propto
  e^{-a V(R) T}.
\end{equation}

A lattice effective potential is therefore defined by the ratio
\begin{equation}
  a V_{\mathrm{eff}}(R,T)
  =
  -
  \ln
  \left[
    \frac{W(R,T+1)}{W(R,T)}
  \right].
  \label{eq:effective-potential-ratio}
\end{equation}
If excited-state contributions are negligible, this approaches the desired
static potential in lattice units:
\begin{equation}
  a V_{\mathrm{eff}}(R,T)
  \xrightarrow[T\to\infty]{}
  a V(R).
\end{equation}

In the numerical analysis, $W(R,T)$ denotes the already normalized and spatially
averaged observable stored as \texttt{W\_temp}. No additional color, volume, or
spatial-direction normalization is applied during the analysis. The main
remaining analysis choices are the thermalization cut, the treatment of
autocorrelations, the temporal range used for the extraction, and the estimation
of statistical uncertainties.

\subsection{Creutz Ratios}
\label{subsec:creutz-ratios}

Creutz ratios are combinations of neighboring Wilson loops designed to cancel
some perimeter and normalization contributions. In the continuum they can be
viewed as the mixed derivative
\begin{equation}
  \chi(r,t)
  =
  -
  \frac{\partial^2}{\partial r\,\partial t}
  \ln W(r,t),
  \label{eq:continuum-creutz-ratio}
\end{equation}
where $W(r,t)$ is a rectangular Wilson loop of physical spatial extent $r$ and
Euclidean time extent $t$. A central finite-difference discretization gives a
lattice Creutz ratio located at half-integer separations,
\begin{equation}
  \chi\left(R+\frac{1}{2},T+\frac{1}{2}\right)
  =
  \ln
  \left[
    \frac{
      W(R+1,T) W(R,T+1)
    }{
      W(R,T) W(R+1,T+1)
    }
  \right],
  \label{eq:creutz-ratio-central-difference}
\end{equation}
or, equivalently after a shift of indices,
\begin{equation}
  \chi(R,T)
  =
  -
  \ln
  \left[
    \frac{
      W(R,T) W(R-1,T-1)
    }{
      W(R,T-1) W(R-1,T)
    }
  \right].
  \label{eq:creutz-ratio}
\end{equation}
The normalization of the Wilson loop, including the factor $1/N_c$, cancels in
this ratio. This is useful for the present $\SU(2)$ analysis because the same
formula can be applied directly to the normalized Wilson-loop output described
above.

The central-difference form approximates the continuum mixed derivative with
cutoff effects of order $\latspacing^2$ when the underlying Wilson loop is
smooth. This makes Creutz ratios particularly useful for continuum studies:
they are finite observables, their leading normalization and perimeter
contributions cancel, and they are closely related to the static force. In the
large-time limit,
\begin{equation}
  \frac{\chi(R,T)}{\latspacing^2}
  \xrightarrow[T\to\infty]{}
  F(r)
\end{equation}
up to the usual distinction between the discrete lattice midpoint and the
physical distance $r$. Diagonal ratios $\chi(R,R)$ are also useful summary
observables because they probe the same spatial and temporal scale and are
straightforward to compare across ensembles \cite{Creutz1980,Risch2024SmearingContinuum}.

If large Wilson loops obey an area law,
\begin{equation}
  W(R,T)
  \sim
  \exp[-\stringtension \latspacing^2 R T
  - \text{perimeter terms}],
\end{equation}
then the Creutz ratio approaches
\begin{equation}
  \chi(R,T)
  \xrightarrow[R,T\to\infty]{}
  \stringtension \latspacing^2 .
\end{equation}
Thus Creutz ratios provide a direct lattice estimator for the string tension in
lattice units.

At finite $R$ and $T$, Creutz ratios contain corrections from excited states,
Coulombic contributions, finite-volume effects, and lattice artifacts. They are
therefore best interpreted either as approximate string-tension estimators at
large loop sizes or as diagnostic observables for scaling comparisons between
ensembles.

For continuum extrapolations the important distinction is between fixed lattice
separation and fixed physical distance. Holding $R$ fixed while taking
$\latspacing\to0$ drives the physical distance $r=R\latspacing$ to zero and
therefore probes the short-distance limit rather than the same physical force.
The continuum comparison should instead use a dimensionless distance such as
\begin{equation}
  \hat r = \frac{r}{\rzero}
\end{equation}
and interpolate the measured Creutz ratios to common $\hat r$ values before
fitting in $\latspacing/\rzero$. This is also where gauge-breaking effects are
expected to be most visible: short-distance Wilson-loop observables are
sensitive to ultraviolet deformations, while large loops suffer from rapidly
increasing statistical noise.

\subsection{Sommer Scale and String Tension}
\label{subsec:sommer-scale-string-tension}

The static potential can also be used to set the scale. A common
phenomenological form is the Cornell-type potential
\begin{equation}
  V(r)
  =
  V_0
  -
  \frac{\alpha}{r}
  +
  \stringtension r,
  \label{eq:cornell-potential}
\end{equation}
where $V_0$ is an additive self-energy, $\alpha$ parametrizes the short-distance
Coulombic part, and $\stringtension$ is the string tension. The corresponding
force is
\begin{equation}
  F(r)
  =
  \frac{\dd V}{\dd r}.
\end{equation}
The additive constant $V_0$ drops out of the force.

The Sommer scale $\rzero$ is defined implicitly by
\begin{equation}
  \rzero^2 F(\rzero) = 1.65 .
  \label{eq:sommer-scale-definition}
\end{equation}
This definition uses the force at intermediate distances, where the potential
can often be determined more accurately than at very large distances. In lattice
units, the force is estimated from finite differences of the potential. If
$\hat V(R)=\latspacing V(R\latspacing)$, then
\begin{equation}
  \hat F\left(R+\frac{1}{2}\right)
  =
  \hat V(R+1) - \hat V(R)
\end{equation}
is a dimensionless estimate of $\latspacing^2 F(r)$ at an intermediate distance.
The Sommer condition can then be written as
\begin{equation}
  \left(\frac{r}{\latspacing}\right)^2
  \hat F\left(\frac{r}{\latspacing}\right)
  =
  1.65 .
  \label{eq:sommer-scale-lattice-units}
\end{equation}
Solving this equation gives $\rzero/\latspacing$.

The ratio $\rzero/\latspacing$ is one of the main scale-setting quantities in
this thesis. It allows data at different $\beta$ values to be compared in
approximately physical units and provides the continuum variable
\begin{equation}
  \frac{\latspacing}{\rzero}
  =
  \left(\frac{\rzero}{\latspacing}\right)^{-1}.
\end{equation}
The continuum limit corresponds to
\begin{equation}
  \frac{\latspacing}{\rzero}
  \to 0.
\end{equation}

The string tension may also be expressed in dimensionless form as
\begin{equation}
  \stringtension \rzero^2 .
\end{equation}
This quantity is useful because it is independent of the arbitrary choice of
lattice units. Agreement of such dimensionless quantities along a sequence of
lattice spacings is a central test of scaling toward a common continuum limit.

\chapter{Small Breaking of Gauge Invariance}
\section{Definition of the Gauge-Breaking Deformation}
Define the modified lattice action or update rule used in the simulations. Specify the breaking parameters, for example $\eps_1$ and $\eps_2$, and explain which gauge transformations are no longer exact symmetries.

\section{Symmetry and Scaling Expectations}
\subsection{Relevant, Marginal, and Irrelevant Breaking Terms}
\subsection{Possible Restoration of Gauge Invariance}
\subsection{Expected Dependence on Breaking Strength and Lattice Spacing}

\section{Working Hypotheses}
List the concrete hypotheses tested numerically:
\begin{enumerate}
  \item Gauge-invariant observables approach the same continuum values for sufficiently small $\eps$.
  \item Gauge-breaking effects scale to zero with a measurable power of $\latspacing$ or $\eps$.
  \item Above a threshold in $\eps$, continuum extrapolations become unstable or incompatible with the gauge-invariant limit.
\end{enumerate}

\chapter{Numerical Setup}
\section{Simulation Parameters}
Describe the gauge group, lattice volumes, coupling range, update algorithm, thermalization strategy, number of sweeps, and stored observables.

\section{Ensemble Organization}
Explain how ensembles are grouped by common physical and algorithmic parameters. State which parameters are varied in the continuum study and which are held fixed.

\section{Thermalization and Autocorrelations}
\subsection{Thermalization Cuts}
\subsection{Integrated Autocorrelation Time}
\subsection{Bootstrap Blocking}

\section{Analysis Pipeline}
Describe the analysis code and reproducibility workflow. Include the path from raw measurements to final observables:
\begin{enumerate}
  \item Load simulation metadata from \texttt{input.yaml}.
  \item Read Wilson-loop and observable measurements.
  \item Apply thermalization cuts.
  \item Estimate uncertainties with bootstrap resampling.
  \item Extract $V(R)$, $\creutz$, $\rzero$, and related quantities.
  \item Aggregate ensembles for continuum and breaking-strength comparisons.
\end{enumerate}

\chapter{Extraction of Physical Observables}
\section{Wilson Loop Analysis}
\subsection{Effective Mass Plateaus}
\subsection{Choice of Fit Windows}
\subsection{Systematic Uncertainty from Fit Ranges}

\section{Static Potential}
\subsection{Potential Extraction}
\subsection{Cornell-Type Fits}
\subsection{Force and Sommer Scale}

\section{Creutz Ratio Analysis}
\subsection{Definition and Interpretation}
\subsection{Large-Time Behavior}
\subsection{Relation to the String Tension}

\section{Scale Setting}
\subsection{Determination of the Sommer Scale in Lattice Units}
\subsection{Alternative Scale Estimates}
\subsection{Consistency Checks}

\chapter{Continuum Extrapolation Method}
\section{Choice of Continuum Variable}
The primary continuum variable is chosen to be the measured dimensionless
lattice spacing
\begin{equation}
  \hat a
  =
  \frac{\latspacing}{\rzero}.
\end{equation}
For observables that are expected to have leading even-power cutoff effects
with the Wilson gauge action, the default horizontal axis is therefore
$\hat a^2=(\latspacing/\rzero)^2$. Using the measured scale rather than the bare
coupling makes the extrapolation closer to a fixed-physics comparison and
reduces the risk of interpreting changes in scale setting as physical changes
in the observable itself.

For Creutz-ratio observables the quantity to extrapolate should be
dimensionless. Convenient choices are
\begin{equation}
  \hat\chi(\hat r)
  =
  \rzero^2
  \frac{\chi(r,t)}{\latspacing^2},
  \qquad
  \hat r = \frac{r}{\rzero},
\end{equation}
or, for large loops where the string-tension interpretation is appropriate,
$\stringtension\rzero^2$. The measured lattice ratio $\chi(R,T)$ is
dimensionless and approaches $\stringtension\latspacing^2$ in the area-law
regime; division by $\latspacing^2$ and multiplication by $\rzero^2$ converts it
to a quantity with a finite physical continuum limit. The same logic applies to
the force extracted from the static potential.

\section{Fixed-Physics Trajectories}
Ensembles at different $\beta$ values are compared along trajectories where the
physical volume and the dimensionless analysis point are kept as fixed as the
available data allow. In practice this means comparing observables at common
$\hat r=r/\rzero$ and using only ensembles whose volumes are large enough that
finite-volume corrections are not expected to dominate the Wilson-loop signal.
When exact matching is not possible, the mismatch is included as a systematic
effect or avoided by restricting the fit range.

The gauge-breaking parameters are treated as additional bare deformation
parameters. If they are held fixed while $\latspacing$ changes, the continuum
fit tests whether their effect vanishes dynamically or survives in the
continuum intercept. If a theoretical scaling prescription for $\eps$ is used,
for example $\eps\propto \latspacing^q$, the fit must instead include this
scaling explicitly. The central diagnostic is whether the broken and unbroken
data approach a common continuum value.

\section{Fit Ansatz}
\subsection{Gauge-Invariant Baseline}
For the gauge-invariant ensembles, a minimal continuum ansatz for a
dimensionless observable $\hat O$ is
\begin{equation}
  \hat O(\hat a)
  =
  O_0
  +
  c_a \hat a^2
  +
  c_{a4} \hat a^4
  +
  \cdots .
  \label{eq:gauge-invariant-continuum-fit}
\end{equation}
The $\hat a^4$ term is included only when the data support it or when it is used
as a fit-systematic variation. The default fit should be stable under removing
the coarsest lattice spacing and under modest changes of the physical-distance
range used for the observable.

\subsection{Combined Lattice-Spacing and Breaking Dependence}
A useful shared-intercept test is
\begin{equation}
  \hat O(\hat a,\eps)
  =
  O_0
  +
  c_a \hat a^2
  +
  c_\eps \eps \hat a^p
  +
  c_{a\eps} \eps \hat a^{p+2}
  +
  \cdots ,
  \label{eq:shared-intercept-breaking-fit}
\end{equation}
with $p>0$. This ansatz encodes the hypothesis that gauge breaking is an
irrelevant cutoff effect: for any fixed small $\eps$ the breaking terms vanish
as $\hat a\to0$, and all data share the same continuum intercept $O_0$.

The competing separate-intercept test is
\begin{equation}
  \hat O(\hat a,\eps)
  =
  O_0
  +
  c_\eps \eps
  +
  c_a \hat a^2
  +
  c_{a\eps} \eps \hat a^2
  +
  \cdots .
  \label{eq:separate-intercept-breaking-fit}
\end{equation}
Here a statistically significant $c_\eps$ indicates that the breaking
deformation changes the continuum value of the observable rather than merely
changing cutoff effects. If the deformation can enter symmetrically in the sign
of $\eps$, a quadratic term $c_{\eps2}\eps^2$ should be included as a systematic
variant.

This structure follows the same logic used in controlled continuum
extrapolations with smearing or flow: one separates changes in the path toward
the continuum from changes in the continuum observable itself
\cite{Risch2024SmearingContinuum}. In that setting, fixed smearing in lattice
units can share the same continuum intercept if it remains in the controlled
expansion regime, while a fixed physical smoothing radius defines a different
continuum observable. In the present gauge-breaking problem there is no
guarantee of a shared intercept, so Eqs.~\eqref{eq:shared-intercept-breaking-fit}
and \eqref{eq:separate-intercept-breaking-fit} are both part of the analysis.

\section{Uncertainty Budget}
\subsection{Statistical Errors}
Statistical errors are obtained from blocked bootstrap samples. The block size
is chosen large enough to reduce residual autocorrelation effects in the saved
Markov-chain history. Derived observables such as Creutz ratios, Sommer scales,
and continuum-fit parameters are recomputed inside the bootstrap whenever
possible, so nonlinear error propagation is captured by the resampling.

\subsection{Fit-Window and Interpolation Systematics}
Wilson-loop observables depend on analysis choices: the temporal range used for
effective-potential plateaus, the loop sizes included in Creutz-ratio summaries,
and the interpolation used to compare fixed physical distances. These choices
are varied to estimate systematic effects. For Creutz ratios, short distances
are most vulnerable to cutoff and breaking artifacts, while long distances are
most vulnerable to statistical noise.

\subsection{Finite-Volume and Continuum-Fit Systematics}
Finite-volume effects are monitored by comparing ensembles with similar
$\hat a$ and $\eps$ but different physical volumes when available. Continuum-fit
systematics are estimated by changing the fit ansatz, removing the coarsest
lattice spacing, and checking whether the extrapolation remains monotonic over
the fitted range. Non-monotonic behavior or strong shifts of the intercept under
reasonable fit variations are treated as evidence that the available lattice
spacings are not yet in a controlled scaling regime.

\chapter{Results}
\section{Thermalization and Data Quality}
Insert representative thermalization histories and autocorrelation estimates.

\section{Gauge-Invariant Baseline}
Present the $\eps=0$ results for $\wilsonloop$, $V(R)$, $\rzero$, and $\stringtension$. This establishes the reference continuum behavior.

\section{Dependence on Gauge-Breaking Strength}
\subsection{Observable Shifts at Fixed Coupling}
\subsection{Scaling of Shifts with Breaking Strength}
\subsection{Thresholds or Nonlinear Behavior}

\section{Continuum Limit at Fixed Breaking}
\subsection{Continuum Extrapolations for Each Breaking Strength}
\subsection{Comparison with the Gauge-Invariant Limit}
\subsection{Stability Under Fit Variations}

\section{Combined Scaling Analysis}
Fit the dependence on both lattice spacing and gauge-breaking parameter. State whether the data support restoration of gauge invariance in the continuum limit.

\section{Finite-Volume and Systematic Checks}
Collect checks that are not part of the primary result but constrain the interpretation.

\chapter{Discussion}
\section{Interpretation of the Continuum Limit}
Discuss whether small gauge breaking behaves like an irrelevant lattice artifact, a relevant deformation, or a source of uncontrolled continuum bias.

\section{Comparison with Theoretical Expectations}
Relate the numerical findings to the hypotheses from Chapter~3.

\section{Limitations}
Discuss limits from statistics, accessible lattice spacings, finite volume, fit-systematic control, and the specific form of gauge breaking studied.

\section{Implications}
Explain what the result means for simulations or formulations where gauge invariance is only approximate.

\chapter{Conclusion and Outlook}
\section{Summary of Findings}
State the central quantitative conclusions.

\section{Answer to the Research Question}
Give the direct answer to whether the continuum limit is robust under the investigated small breaking of gauge invariance.

\section{Future Work}
Possible extensions:
\begin{itemize}
  \item Study additional forms of gauge-breaking terms.
  \item Increase statistics and extend to finer lattice spacings.
  \item Add gauge-variant diagnostics to directly measure symmetry restoration.
  \item Compare different gauge groups or dimensions.
\end{itemize}

\appendix

\chapter{Simulation Tables}
Provide complete tables of ensembles, including $\beta$, lattice size, breaking parameters, number of sweeps, thermalization cuts, and measured scales.

\chapter{Additional Plots}
Include supplementary Wilson-loop plateaus, bootstrap block-size scans, potential fits, and continuum-fit variants.

\chapter{Analysis Code and Reproducibility}
Document the exact command-line workflow, cache version, software environment, and data-selection rules used for the final analysis.

\chapter{Derivations}
Collect longer derivations that would interrupt the main text, such as Creutz-ratio identities or continuum-scaling arguments.

\bibliographystyle{unsrt}
\bibliography{references}

\end{document}
